\begin{abstract}
	Small bioactivity data sets make it easy to overstate the benefit of complex models, through data leakage and through model selection that is rarely evaluated on unseen data. We re-examined this for inhibitors of \textit{Toxoplasma gondii} dihydrofolate reductase (TgDHFR, $N = \NCompounds$ curated compounds), comparing classical stacking ensembles on four molecular representations with a Directed Message Passing Neural Network (D-MPNN) on the same \TotalOuterFolds\ folds of a repeated nested cross-validation. Ensembles on Morgan fingerprints reached $R^2 = \MorganRTwoMean \pm \MorganRTwoStd$; adding 3D descriptors changed nothing detectable ($\Delta R^2 = \NBdiffRTwoTwoDFpVsTwoDThreeDFp$). The D-MPNN reached $R^2 = \ChemPropDMPNNFoldRTwoMean$ with default settings and $\GnnTunedRTwoMean$ after a nested hyperparameter search, not significantly different from any classical representation after multiple-comparison correction (adjusted $p \geq \GnnTunedVsMinPRTwo$); equivalence with Morgan fingerprints held for the MAE but not for $R^2$. Feature selection and Gaussian-noise augmentation, credited with large gains in the original study, did not improve a classical learner, so those gains were not reproduced. With scaffold-grouped splits, $R^2$ fell by about 0.2 for both families. The classical benchmark was not cheaper to run than the network's. Applying the Morgan ensemble to an FDA-approved library (\ScreenNLibrary\ compounds) and docking \ScreenNCandidates\ candidates gave \AdmetNConsensus\ consensus hits, \AdmetNConsensusSurvivors\ of which passed in silico ADMET filters; potency--docking agreement was weak ($\rho = \DockSpearmanRho$) and selectivity over human DHFR was not assessed. For this target, a fixed-fingerprint ensemble was sufficient; we make no claim beyond it.
\end{abstract}

\textbf{Keywords:} 
Quantitative Structure-Activity Relationship, Low-Data Regimes, Nested Cross-Validation, Data Leakage, Graph Neural Networks, Molecular Docking